\section{Thông tin nhóm}

Báo cáo được thực hiện bởi nhóm 13 của lớp 23\_24 môn Nhập môn học máy.

\subsection{Thông tin thành viên}
Thông tin các thành viên nhóm được trình bày ở
\tablename~\ref{tab:info-members}.

\begin{longtblr}[
  caption = {Thông tin thành viên nhóm},
  label = {tab:info-members},
  note{$\dag$} = {Nhóm trưởng.}
]{
  width=\textwidth,
  colspec={lll},
  rowhead = 1,
  hlines,
  vlines
}
  Tên thành viên & Mã số sinh viên\\
  Hoàng Ngọc Phú\TblrNote{$\dag$} & 23120010\\
  Hoàng Ngọc Quí & 23120077\\
  Nguyễn Duy Bảo & 23120113 \\
\end{longtblr}

\subsection{Phân chia công việc và mức độ hoàn thành}

\tablename~\ref{tab:task-members} trình bày các nhiệm vụ và thành viên thực
hiện. Mức độ hoàn thành được tính riêng cho từng thành viên thay vì so với tỉ
lệ hoàn thành trong từng công việc---tức là ai hoàn thành công việc của mình
đều được đánh giá 100\%.

\begin{longtblr}[
  caption = {Nhiệm vụ và mức độ hoàn thành yêu cầu được giao của mỗi thành
  viên},
  label = {tab:task-members}
]{
  width=\textwidth,
  colspec={Q[3in,h,l]*{3}{X[c,m]}},
  cell{1-2}{1} = {font=\bfseries, bg=gray!15},
  cell{1}{2-Z} = {font=\bfseries, bg=gray!15},
  rowhead = 2,
  hlines,
  vlines
}
  \SetCell[r=2]{c,m} Nhiệm vụ & \SetCell[c=3]{c} Thành viên & \\
   & H.\ N.\ Phú & H.\ N.\ Quí & N.\ D.\ Bảo\\
  \SetCell[c=4]{l,font=\bfseries} Bài toán hồi quy & \\
  Chọn bộ dữ liệu & 100\% & 100\% & 100\% \\
  Phân tích khám phá và tiền xử lí & 100\% \\
  Hồi quy tuyến tính & 100\% \\
  Hồi quy có Regularization và Lựa chọn Đặc trưng & & 100\% \\
  Mô hình với Hàm Cơ sở Phi Tuyến và Ablation Study & & & 100\% \\
  Hồi quy Bayesian đầy đủ + vùng bất định (nâng cao) & & & 100\% \\
  Evidence Maximization (nâng cao) & & & 100\% \\
  Kernel Ridge Regression (nâng cao) & & 100\% \\
  Gaussian Process Regression (nâng cao) & 100\% \\
  Robust Regression (nâng cao) & 100\% \\
  Phân tích Bias–Variance thực nghiệm (nâng cao) & & 100\% \\
  Đánh giá mô hình & 100\% & 100\% & 100\% \\
  Phân tích và thảo luận & 100\% & 100\% & 100\% \\
  \SetCell[c=4]{l,font=\bfseries} Bài toán phân lớp \\
  Chọn bộ dữ liệu & 100\% & 100\% & 100\% \\
  Phân tích khám phá và tiền xử lí & 100\% \\
  Logistic Regression & 100\% \\
  Linear Discriminant Analysis--LDA và QDA & & 100\% \\
  Perceptron và Logistic Regression có Regularization & & & 100\% \\
  Mô hình Probit (nâng cao) & 100\% \\
  Xấp xỉ Laplace (nâng cao) & 100\% \\
  Kernel Logistic Regression (nâng cao) & & 100\% \\
  Gaussian Naive Bayes vs.\ LDA (nâng cao) & & 100\% \\
  Phân tích VC dimension (nâng cao) & & & 100\% \\
  Đánh giá mô hình & 100\% & 100\% & 100\% \\
  Phân tích và thảo luận & 100\% & 100\% & 100\% \\
  Kết nối hồi quy và phân lớp & & 100\% & 100\% \\
  Bảng so sánh toàn diện & 100\% & 100\% & 100\% \\
  Phân tích mức độ nhạy cảm (sensitivity) và tính bền vững (robustness) & 100\%
  \\
  \SetCell[c=4]{l,font=\bfseries} Khác \\
  Cài đặt & 100\% & 100\% & 100\% \\
  Chạy thí nghiệm & 100\% & 100\% & 100\% \\
  Viết báo cáo & 100\% & 100\% & 100\% \\
\end{longtblr}

\subsection{Mức độ hoàn thành đồ án}

\tablename~\ref{tab:requirement-project} trình bày từng yêu cầu của đồ án và
mức độ hoàn thành.

\begin{longtblr}[
  caption = {Yêu cầu và mức độ hoàn thành từng yêu cầu},
  label = {tab:requirement-project},
  note{$\dag$} = {Chưa thực hiện ở các phần nâng cao.}
]{
  width=\textwidth,
  colspec={Q[3in,h,l]Q[1in,c,m]},
  row{1} = {font=\bfseries, bg=gray!15},
  rowhead = 1,
  hlines,
  vlines
}
  \SetRow[]{c} Yêu cầu & Mức độ hoàn thành \\
  \SetCell[c=2]{l,font=\bfseries} Bài toán hồi quy \\
  Cơ sở lý thuyết: Normal Eq., Gauss–Markov, Bias–Variance, Bayesian, Evidence
  & 100\% \\
  EDA, kiểm định Breusch–Pagan, WLS nếu cần, tiền xử lý đầy đủ & 100\% \\
  Linear Regression: Normal Equations + Mini-batch GD + learning
  rate schedule & 100\% \\
  Elastic Net + Lasso path + Forward/Backward feature selection & 100\% \\
  Mô hình phi tuyến (3 loại) + validation curve + ablation study & 100\% \\
  Đánh giá: $k$-fold CV $k = 10$, kiểm định thống kê, calibration & 100\% \\
  Hồi quy Bayesian (nâng cao) & 100\% \\
  Evidence Maximization (nâng cao) & 100\% \\
  Kernel Ridge (nâng cao) & 100\% \\
  GP Regression (nâng cao) & 100\% \\
  Robust Regression (nâng cao) & 100\% \\
  Bias–-Var bootstrap & 100\% \\
  \SetCell[c=2]{l,font=\bfseries} Bài toán phân lớp \\

  Lý thuyết: gradient/Hessian LR, IRLS, Perceptron, Kernel LR, Laplace, Probit
  & 100\% \\
  EDA, phân tích mất cân bằng, tiền xử lý có chiều sâu & 100\% \\
  Logistic Regression: GD + Newton–Raphson + so sánh 3 chiến lược đa lớp &
  100\% \\
  LDA + QDA: Fisher ratio, so sánh lý thuyết, quy chiếu 2D & 100\% \\
  Perceptron + L1/L2 LR + class-weighted loss + stratified CV & 100\% \\
  Đánh giá: CV, McNemar, PR curve, calibration, phân tích lỗi sâu & 100\% \\
  Probit (nâng cao) & 100\% \\
  Laplace (nâng cao) & 100\% \\
  Kernel LR (nâng cao) & 100\% \\
  GNB vs LDA (nâng cao) & 100\% \\
  VC dimension (nâng cao) & 100\% \\
  \SetCell[c=2]{l,font=\bfseries} Kết nối Hồi quy và Phân lớp & 100\% \\
  Hồi quy logistic như hồi quy & 100\% \\
  Generalized Linear Models (GLM) & 100\% \\
  Exponential family & 100\% \\
  Bảng so sánh toàn diện & 100\% \\
  \SetCell[c=2]{l,font=\bfseries} Phân tích mức độ nhạy cảm (sensitivity) và
  tính bền vững (robustness) \\
  Sensitivity analysis & 100\% \\
  Noise injection & 100\% \\
  Feature corruption & 100\% \\
  Phân tích khả năng hội tụ (convergence) & 100\% \\
  \SetCell[c=2]{l,font=\bfseries} Khả năng tái hiện lại thí nghiệm \\
  Cố định tất cả random seed & 100\% \\
  Ghi log đầy đủ: hyperparameter, train/val/test split indices, metric từng
  fold, $\ldots\,$ & 80\%\TblrNote{$\dag$} \\
  Báo cáo thời gian thực thi/ huấn luyện/ kiểm tra và bộ nhớ tiêu tốn của từng
  thuật toán & 100\% \\
  Cung cấp file \texttt{environment.yml} (conda) hoặc \texttt{requirements.txt}
  với version cụ thể & 100\% \\
  Kết quả phải tái hiện được hoàn toàn khi chạy lại notebook từ đầu & 100\% \\
  \SetCell[c=2]{l,font=\bfseries} Yêu cầu về Code \\
  Tất cả code phải chạy lại được & 100\% \\
  Phần cài đặt từ đầu phải rõ ràng, tách biệt với phần sử dụng thư viện kiểm
  chứng & 100\% \\
  Phiên bản Python và thư viện phải được ghi rõ trong file
  \texttt{requirements.txt} & 100\% \\
  Đặt seed ngẫu nhiên để đảm bảo kết quả tái hiện được & 100\% \\
  \SetCell[c=2]{l,font=\bfseries} Yêu cầu về Báo cáo \\
  Báo cáo được nộp dưới dạng PDF, không giới hạn độ dài trang & 100\% \\
  Danh mục tài liệu tham khảo theo chuẩn IEEE hoặc APA rõ ràng & 100\% \\
  Khuyến nghị sử dụng \LaTeX. Nếu sử dụng \LaTeX, yêu cầu nhóm nộp cả mã nguồn
  \texttt{.tex} & 100\% \\
  Mỗi thành viên trong nhóm phải có mục đóng góp rõ ràng & 100\% \\
\end{longtblr}
